% Part: normal-modal-logic
% Chapter: syntax-and-semantics
% Section: language-modal-logic

\documentclass[../../../include/open-logic-section]{subfiles}

\begin{document}

\olfileid[pt]{nml}{syn}{lan}

\olsection{A linguagem da lógica modal básica}

\begin{defn}
A linguagem básica da lógica modal contém:
\begin{enumerate}
  \tagitem{prvFalse}{A constante proposicional para !!{falsity}~$\lfalse$.}{}
  \tagitem{prvTrue}{A constante proposicional para !!{truth}~$\ltrue$.}{}
  \item Um conjunto !!{denumerable}s de !!{propositional variable}s: $\Obj
    p_0$, $\Obj p_1$, $\Obj p_2$, \dots
  \item Os conectivos proposicionais: \startycommalist
  \iftag{prvNot}{\ycomma $\lnot$ (negação)}{}%
  \iftag{prvAnd}{\ycomma $\land$ (conjunção)}{}%
  \iftag{prvOr}{\ycomma $\lor$ (disjunção)}{}%
  \iftag{prvIf}{\ycomma $\lif$ (!!{conditional})}{}%
  \iftag{prvIff}{\ycomma $\liff$ (!!{biconditional})}{}.
  \tagitem{prvBox}{O operador modal $\Box$.}{}
  \tagitem{prvDiamond}{O operador modal $\Diamond$.}{}
\end{enumerate}
\end{defn}

\begin{defn}
As \emph{!!^{formula}s} da linguagem modal básica são definidas
indutivamente como segue:
\begin{enumerate}
\tagitem{prvFalse}{$\lfalse$ é uma !!{formula} atômica.}{}

\tagitem{prvTrue}{$\ltrue$ é uma !!{formula} atômica.}{}

\item Toda !!{propositional variable} $\Obj p_i$ é uma !!{formula} (atômica).

\tagitem{prvNot}{Se $!A$ é !!a{formula}, então $\lnot !A$ é
  !!a{formula}.}{}

\tagitem{prvAnd}{Se $!A$ e $!B$ são !!{formula}s, então $(!A \land
  !B)$ é !!a{formula}.}{}

\tagitem{prvOr}{Se $!A$ e $!B$ são !!{formula}s, então $(!A \lor !B)$
  é !!a{formula}.}{}

\tagitem{prvIf}{Se $!A$ e $!B$ são !!{formula}s, então $(!A \lif !B)$
  é !!a{formula}.}{}

\tagitem{prvIff}{Se $!A$ e $!B$ são !!{formula}s, então $(!A \liff !B)$
  é !!a{formula}.}{}

\tagitem{prvBox}{Se $!A$ é !!a{formula}, então $\Box !A$ é
  !!a{formula}.}{}

\tagitem{prvDiamond}{Se $!A$ é !!a{formula}, então $\Diamond !A$ é
  !!a{formula}.}{}

\tagitem{limitClause}{Nada mais é !!a{formula}.}{}
\end{enumerate}
\end{defn}

\begin{tagblock}{defNot,defOr,defAnd,defIf,defIff,defTrue,defFalse,defBox,defDiamond}
\begin{defn}
As fórmulas construídas com os operadores definidos devem ser entendidas
como segue:

\begin{tagenumerate}{defTrue,defFalse,defNot,defOr,defAnd,defIf,defIff,defBox,defDiamond}

\tagitem{defTrue}{$\ltrue$ abrevia
  \iftag{prvFalse}{$\lnot\lfalse$}{$(!A \lor \lnot !A)$ para alguma
    !!{formula} atômica fixa~$!A$}.}{}

\tagitem{defFalse}{$\lfalse$ abrevia
  \iftag{prvTrue}{$\lnot\ltrue$}{$(!A \land \lnot !A)$ para alguma
    !!{formula} atômica fixa~$!A$}.}{}

\tagitem{defNot}{$\lnot !A$ abrevia $!A \lif \lfalse$.}{}

\tagitem{defOr}{$!A \lor !B$ abrevia
  \iftag{prvAnd}{$\lnot(\lnot !A \land \lnot !B)$}{$\lnot !A \lif
    !B$}.}{}

\tagitem{defAnd}{$!A \land !B$ abrevia
  \iftag{prvOr}{$\lnot(\lnot !A \lor \lnot !B)$}{$\lnot (!A \lif
    \lnot !B)$}.}{}

\tagitem{defIf}{$!A \lif !B$ abrevia
  \iftag{prvOr}{$\lnot !A \lor !B)$}{$\lnot (!A \land \lnot !B)$}.}{}

\tagitem{defIff}{$!A \liff !B$ abrevia $(!A \lif !B) \land (!B
  \lif !A)$.}{}

\tagitem{defBox}{$\Box !A$ abrevia $\lnot\Diamond\lnot !A$. }{}

\tagitem{defDiamond}{$\Diamond !A$ abrevia $\lnot\Box\lnot !A$.}{}
\end{tagenumerate}
\end{defn}
\end{tagblock}

Se !!a{formula}~$!A$ não contém $\Box$ nem $\Diamond$, dizemos que ela
é \emph{livre de operadores modais}.

\end{document}
